\begin{problem}{Sieve of Sundaram}
    In 1934, S. P. Sundaram introduced a method to generate all odd primes by removing numbers of the form \(i + j + 2ij\) from the integers. Specifically, for \(i, j \geq 1\), mark every integer of the form \(i + j + 2ij\). The remaining unmarked numbers \(k\), when doubled and incremented, give us odd primes \(2k+1\).

    Let \(F(N)\) be the number of integers in \(\{1, 2, \ldots, N\}\) that are never marked, that is, cannot be written as \(i + j + 2ij\) for any \(i, j \geq 1\). Compute \(F(10^6)\).

    
\end{problem}
